% Gesamtarchitektur der Umgebung: Bedienwege, Funktionsknoten, Herstellerschnittstelle, Simulation oder Roboter.
% Eigene Darstellung (Art E, Kap. 6.2.1). Keine Quelle in der Caption.
% Basis: G1@25d5c1a g1pilot/docs/02-T-Architektur.md (Gesamtbild, Laufzeit-Topologie, Zwei Kommunikationsebenen,
%   Sim <-> Real: Ein Schalter, Joint-Layout/Zuständigkeit der Gelenke), docs/05-T-Nodes-und-Topics.md,
%   docs/10-T-Demo-GUI.md (Grundsatz), docs/30-T-Stehen-und-Gehen.md (Warum zwei Implementierungen),
%   docs/21-T-Arm-API.md (Arm-API als eigener Knoten), docs/50-A-Echter-Roboter.md Abschnitt E und
%   docs/22-T-Inspire-Haende.md (Hände real getrennter Anschluss, Sim über Simulation).
% Vereinfachungen: Navigation (nur Simulation) und Hand-Profile nicht gezeichnet; die Bedienoberfläche läuft
%   technisch im selben Container wie die Funktionsknoten, steht hier als Bedienweg oben.
% Stil: ba-grafik.tex. Einbinden: % Gesamtarchitektur der Umgebung: Bedienwege, Funktionsknoten, Herstellerschnittstelle, Simulation oder Roboter.
% Eigene Darstellung (Art E, Kap. 6.2.1). Keine Quelle in der Caption.
% Basis: G1@25d5c1a g1pilot/docs/02-T-Architektur.md (Gesamtbild, Laufzeit-Topologie, Zwei Kommunikationsebenen,
%   Sim <-> Real: Ein Schalter, Joint-Layout/Zuständigkeit der Gelenke), docs/05-T-Nodes-und-Topics.md,
%   docs/10-T-Demo-GUI.md (Grundsatz), docs/30-T-Stehen-und-Gehen.md (Warum zwei Implementierungen),
%   docs/21-T-Arm-API.md (Arm-API als eigener Knoten), docs/50-A-Echter-Roboter.md Abschnitt E und
%   docs/22-T-Inspire-Haende.md (Hände real getrennter Anschluss, Sim über Simulation).
% Vereinfachungen: Navigation (nur Simulation) und Hand-Profile nicht gezeichnet; die Bedienoberfläche läuft
%   technisch im selben Container wie die Funktionsknoten, steht hier als Bedienweg oben.
% Stil: ba-grafik.tex. Einbinden: % Gesamtarchitektur der Umgebung: Bedienwege, Funktionsknoten, Herstellerschnittstelle, Simulation oder Roboter.
% Eigene Darstellung (Art E, Kap. 6.2.1). Keine Quelle in der Caption.
% Basis: G1@25d5c1a g1pilot/docs/02-T-Architektur.md (Gesamtbild, Laufzeit-Topologie, Zwei Kommunikationsebenen,
%   Sim <-> Real: Ein Schalter, Joint-Layout/Zuständigkeit der Gelenke), docs/05-T-Nodes-und-Topics.md,
%   docs/10-T-Demo-GUI.md (Grundsatz), docs/30-T-Stehen-und-Gehen.md (Warum zwei Implementierungen),
%   docs/21-T-Arm-API.md (Arm-API als eigener Knoten), docs/50-A-Echter-Roboter.md Abschnitt E und
%   docs/22-T-Inspire-Haende.md (Hände real getrennter Anschluss, Sim über Simulation).
% Vereinfachungen: Navigation (nur Simulation) und Hand-Profile nicht gezeichnet; die Bedienoberfläche läuft
%   technisch im selben Container wie die Funktionsknoten, steht hier als Bedienweg oben.
% Stil: ba-grafik.tex. Einbinden: % Gesamtarchitektur der Umgebung: Bedienwege, Funktionsknoten, Herstellerschnittstelle, Simulation oder Roboter.
% Eigene Darstellung (Art E, Kap. 6.2.1). Keine Quelle in der Caption.
% Basis: G1@25d5c1a g1pilot/docs/02-T-Architektur.md (Gesamtbild, Laufzeit-Topologie, Zwei Kommunikationsebenen,
%   Sim <-> Real: Ein Schalter, Joint-Layout/Zuständigkeit der Gelenke), docs/05-T-Nodes-und-Topics.md,
%   docs/10-T-Demo-GUI.md (Grundsatz), docs/30-T-Stehen-und-Gehen.md (Warum zwei Implementierungen),
%   docs/21-T-Arm-API.md (Arm-API als eigener Knoten), docs/50-A-Echter-Roboter.md Abschnitt E und
%   docs/22-T-Inspire-Haende.md (Hände real getrennter Anschluss, Sim über Simulation).
% Vereinfachungen: Navigation (nur Simulation) und Hand-Profile nicht gezeichnet; die Bedienoberfläche läuft
%   technisch im selben Container wie die Funktionsknoten, steht hier als Bedienweg oben.
% Stil: ba-grafik.tex. Einbinden: \input{Bilder/tikz/architektur-gesamt} in einer figure-Umgebung. Nicht skalieren.
\begin{tikzpicture}[ba]
  \tikzset{
    bw/.style={ba block, text width=31mm, minimum height=10mm},
    fk/.style={ba block hervor, minimum height=17mm, minimum width=0pt},
    sub/.style={ba block, font=\footnotesize, minimum height=10mm, minimum width=0pt, inner sep=2pt},
    ziel/.style={ba pfeil},
    zust/.style={ba pfeil, dash pattern=on 3pt off 2pt},
    klein/.style={font=\footnotesize, text=baGrau},
  }
  % ---------------- Ebene 1: Bedienwege ----------------
  \node[klein, anchor=south west] at (1.3,11.45) {Bedienwege (rechnen keine Bewegung, senden Ziele und lesen Zustände)};
  \node[bw] (B1) at (3.0,10.65) {Bedien-\\oberfläche};
  \node[bw] (B2) at (6.6,10.65) {Gamepad};
  \node[bw] (B3) at (10.2,10.65) {Smartphone};
  \node[bw] (B4) at (13.8,10.65) {Externe\\Programme};
  % Sammelleitung ROS 2
  \draw[ba normal] (3.0,9.95) -- (13.8,9.95);
  \foreach \x in {3.0,6.6,10.2,13.8} {\draw[ba normal] (\x,10.15) -- (\x,9.95);}
  \node[klein, fill=white, inner sep=1.5pt] at (4.3,9.95) {ROS 2};
  % ---------------- Container Umgebung ----------------
  \draw[ba hilfslinie] (1.2,5.25) rectangle (15.6,9.35);
  \node[klein, anchor=south west] at (1.3,5.3) {Container Umgebung};
  % Ebene 2: Funktionsknoten
  \node[fk, text width=22mm] (RZ) at (2.5,7.4) {Roboter-\\zustand};
  \node[fk, text width=28mm] (AR) at (5.9,7.4) {Armregelung\\mit Planung};
  \node[fk, minimum width=51mm, minimum height=24mm, anchor=center] (FB) at (10.55,7.4) {};
  \node[font=\small, anchor=north] at (10.55,8.55) {Fortbewegung};
  \node[sub, text width=21mm] (ER) at (9.2,7.15) {Ersatzregler\\(Simulation)};
  \node[sub, text width=21mm] (OA) at (11.9,7.15) {Anbindung Onboard-Regler (Realbetrieb)};
  \node[fk, text width=19mm] (HA) at (14.35,7.4) {Hände};
  % Pfeile Bedienwege -> Funktionsknoten (Ziele nach unten, Zustände nach oben)
  \draw[ziel] (5.6,9.95) -- (5.6,8.6);
  \draw[zust] (6.2,8.6) -- (6.2,9.95);
  \draw[ziel] (8.6,9.95) -- (8.6,8.6);
  \draw[zust] (12.5,8.6) -- (12.5,9.95);
  \draw[ziel] (14.05,9.95) -- (14.05,8.6);
  \draw[zust] (14.65,8.6) -- (14.65,9.95);
  \draw[zust] (2.5,8.25) -- (2.5,9.95);
  % Schutzmechanismen
  \node[ba block hervor, minimum height=6mm, minimum width=108mm, font=\footnotesize] (SM) at (10.1,5.95)
    {Schutzmechanismen: begrenzen \textperiodcentered{} vermeiden \textperiodcentered{} überwachen \textperiodcentered{} anhalten};
  \draw[ba duenn, color=baGruen] (5.9,6.55) -- (5.9,6.25);
  \draw[ba duenn, color=baGruen] (10.55,6.55) -- (10.55,6.25);
  \draw[ba duenn, color=baGruen] (14.35,6.55) -- (14.35,6.25);
  % ---------------- Ebene 3: Herstellerschnittstelle ----------------
  \node[ba block, fill=baHellgrau!40, minimum height=8mm, minimum width=143mm] (DDS) at (8.4,4.2)
    {Herstellerschnittstelle (Unitree-DDS)};
  % Befehle nach unten, Zustände nach oben zwischen Funktionsknoten und Schnittstelle
  \draw[zust] (2.5,4.6) -- (2.5,6.55);
  \draw[ziel] (5.6,5.65) -- (5.6,4.6);
  \draw[zust] (6.2,4.6) -- (6.2,5.65);
  \draw[ziel] (9.5,5.65) -- (9.5,4.6);
  \draw[zust] (10.1,4.6) -- (10.1,5.65);
  \draw[ziel] (11.5,5.65) -- (11.5,4.6);
  \draw[zust] (12.1,4.6) -- (12.1,5.65);
  \draw[ziel] (14.05,5.65) -- (14.05,4.6);
  \draw[zust] (14.65,4.6) -- (14.65,5.65);
  % ---------------- Umschalter ----------------
  \node[draw=baAkzent, ba normal, fill=white, rounded corners=1.5pt, text=black, align=center,
        minimum height=9mm, inner sep=3pt] (SW) at (8.4,2.75) {Umschalter\\Simulation / Realbetrieb};
  \draw[baAkzent, ba normal] (8.4,3.8) -- (SW.north);
  % ---------------- Ebene 4: Simulation oder Roboter ----------------
  \draw[ba gruppe] (1.3,-1.1) rectangle (7.0,1.5);
  \node[klein, anchor=north west] at (1.4,1.5) {Container Simulation};
  \node[ba block, minimum width=0pt, text width=22mm] (SIM) at (2.9,-0.1) {MuJoCo-Simulation};
  \node[ba block, minimum width=0pt, text width=22mm] (SH) at (5.5,-0.1) {Simulierte Hände};
  \draw[ba gruppe] (9.8,-1.1) rectangle (15.5,1.5);
  \node[klein, anchor=north west] at (9.9,1.5) {Realbetrieb};
  \node[ba block, minimum width=0pt, text width=22mm] (RB) at (11.4,-0.1) {Unitree G1 mit Onboard-Regler};
  \node[ba block, minimum width=0pt, text width=22mm] (RH) at (14.0,-0.1) {Inspire-Hände};
  \draw[baAkzent, ba normal, <->] (SW.west) -| (4.15,1.5);
  \draw[baAkzent, ba normal, <->] (SW.east) -| (12.65,1.5);
  \node[font=\footnotesize, anchor=east, align=right] at (4.05,2.1) {Domain 1,\\Schnittstelle Loopback};
  \node[font=\footnotesize, anchor=west, align=left] at (12.75,2.1) {Domain 0,\\Netzwerkkarte};
  % Hände real: eigener Anschluss neben der Herstellerschnittstelle
  \draw[ba normal, dash pattern=on 1pt off 2pt, <->] (HA.east) -- (15.85,7.4) -- (15.85,-0.1) -- (RH.east);
  \node[klein, anchor=east, align=right, fill=white, inner sep=1pt] at (15.8,3.2) {eigener Anschluss,\\nur Realbetrieb};
  % Legende
  \draw[ziel] (1.3,-1.75) -- (2.2,-1.75);
  \node[klein, anchor=west] at (2.3,-1.75) {Ziele und Befehle};
  \draw[zust] (6.0,-1.75) -- (6.9,-1.75);
  \node[klein, anchor=west] at (7.0,-1.75) {Zustände und Meldungen};
\end{tikzpicture}
 in einer figure-Umgebung. Nicht skalieren.
\begin{tikzpicture}[ba]
  \tikzset{
    bw/.style={ba block, text width=31mm, minimum height=10mm},
    fk/.style={ba block hervor, minimum height=17mm, minimum width=0pt},
    sub/.style={ba block, font=\footnotesize, minimum height=10mm, minimum width=0pt, inner sep=2pt},
    ziel/.style={ba pfeil},
    zust/.style={ba pfeil, dash pattern=on 3pt off 2pt},
    klein/.style={font=\footnotesize, text=baGrau},
  }
  % ---------------- Ebene 1: Bedienwege ----------------
  \node[klein, anchor=south west] at (1.3,11.45) {Bedienwege (rechnen keine Bewegung, senden Ziele und lesen Zustände)};
  \node[bw] (B1) at (3.0,10.65) {Bedien-\\oberfläche};
  \node[bw] (B2) at (6.6,10.65) {Gamepad};
  \node[bw] (B3) at (10.2,10.65) {Smartphone};
  \node[bw] (B4) at (13.8,10.65) {Externe\\Programme};
  % Sammelleitung ROS 2
  \draw[ba normal] (3.0,9.95) -- (13.8,9.95);
  \foreach \x in {3.0,6.6,10.2,13.8} {\draw[ba normal] (\x,10.15) -- (\x,9.95);}
  \node[klein, fill=white, inner sep=1.5pt] at (4.3,9.95) {ROS 2};
  % ---------------- Container Umgebung ----------------
  \draw[ba hilfslinie] (1.2,5.25) rectangle (15.6,9.35);
  \node[klein, anchor=south west] at (1.3,5.3) {Container Umgebung};
  % Ebene 2: Funktionsknoten
  \node[fk, text width=22mm] (RZ) at (2.5,7.4) {Roboter-\\zustand};
  \node[fk, text width=28mm] (AR) at (5.9,7.4) {Armregelung\\mit Planung};
  \node[fk, minimum width=51mm, minimum height=24mm, anchor=center] (FB) at (10.55,7.4) {};
  \node[font=\small, anchor=north] at (10.55,8.55) {Fortbewegung};
  \node[sub, text width=21mm] (ER) at (9.2,7.15) {Ersatzregler\\(Simulation)};
  \node[sub, text width=21mm] (OA) at (11.9,7.15) {Anbindung Onboard-Regler (Realbetrieb)};
  \node[fk, text width=19mm] (HA) at (14.35,7.4) {Hände};
  % Pfeile Bedienwege -> Funktionsknoten (Ziele nach unten, Zustände nach oben)
  \draw[ziel] (5.6,9.95) -- (5.6,8.6);
  \draw[zust] (6.2,8.6) -- (6.2,9.95);
  \draw[ziel] (8.6,9.95) -- (8.6,8.6);
  \draw[zust] (12.5,8.6) -- (12.5,9.95);
  \draw[ziel] (14.05,9.95) -- (14.05,8.6);
  \draw[zust] (14.65,8.6) -- (14.65,9.95);
  \draw[zust] (2.5,8.25) -- (2.5,9.95);
  % Schutzmechanismen
  \node[ba block hervor, minimum height=6mm, minimum width=108mm, font=\footnotesize] (SM) at (10.1,5.95)
    {Schutzmechanismen: begrenzen \textperiodcentered{} vermeiden \textperiodcentered{} überwachen \textperiodcentered{} anhalten};
  \draw[ba duenn, color=baGruen] (5.9,6.55) -- (5.9,6.25);
  \draw[ba duenn, color=baGruen] (10.55,6.55) -- (10.55,6.25);
  \draw[ba duenn, color=baGruen] (14.35,6.55) -- (14.35,6.25);
  % ---------------- Ebene 3: Herstellerschnittstelle ----------------
  \node[ba block, fill=baHellgrau!40, minimum height=8mm, minimum width=143mm] (DDS) at (8.4,4.2)
    {Herstellerschnittstelle (Unitree-DDS)};
  % Befehle nach unten, Zustände nach oben zwischen Funktionsknoten und Schnittstelle
  \draw[zust] (2.5,4.6) -- (2.5,6.55);
  \draw[ziel] (5.6,5.65) -- (5.6,4.6);
  \draw[zust] (6.2,4.6) -- (6.2,5.65);
  \draw[ziel] (9.5,5.65) -- (9.5,4.6);
  \draw[zust] (10.1,4.6) -- (10.1,5.65);
  \draw[ziel] (11.5,5.65) -- (11.5,4.6);
  \draw[zust] (12.1,4.6) -- (12.1,5.65);
  \draw[ziel] (14.05,5.65) -- (14.05,4.6);
  \draw[zust] (14.65,4.6) -- (14.65,5.65);
  % ---------------- Umschalter ----------------
  \node[draw=baAkzent, ba normal, fill=white, rounded corners=1.5pt, text=black, align=center,
        minimum height=9mm, inner sep=3pt] (SW) at (8.4,2.75) {Umschalter\\Simulation / Realbetrieb};
  \draw[baAkzent, ba normal] (8.4,3.8) -- (SW.north);
  % ---------------- Ebene 4: Simulation oder Roboter ----------------
  \draw[ba gruppe] (1.3,-1.1) rectangle (7.0,1.5);
  \node[klein, anchor=north west] at (1.4,1.5) {Container Simulation};
  \node[ba block, minimum width=0pt, text width=22mm] (SIM) at (2.9,-0.1) {MuJoCo-Simulation};
  \node[ba block, minimum width=0pt, text width=22mm] (SH) at (5.5,-0.1) {Simulierte Hände};
  \draw[ba gruppe] (9.8,-1.1) rectangle (15.5,1.5);
  \node[klein, anchor=north west] at (9.9,1.5) {Realbetrieb};
  \node[ba block, minimum width=0pt, text width=22mm] (RB) at (11.4,-0.1) {Unitree G1 mit Onboard-Regler};
  \node[ba block, minimum width=0pt, text width=22mm] (RH) at (14.0,-0.1) {Inspire-Hände};
  \draw[baAkzent, ba normal, <->] (SW.west) -| (4.15,1.5);
  \draw[baAkzent, ba normal, <->] (SW.east) -| (12.65,1.5);
  \node[font=\footnotesize, anchor=east, align=right] at (4.05,2.1) {Domain 1,\\Schnittstelle Loopback};
  \node[font=\footnotesize, anchor=west, align=left] at (12.75,2.1) {Domain 0,\\Netzwerkkarte};
  % Hände real: eigener Anschluss neben der Herstellerschnittstelle
  \draw[ba normal, dash pattern=on 1pt off 2pt, <->] (HA.east) -- (15.85,7.4) -- (15.85,-0.1) -- (RH.east);
  \node[klein, anchor=east, align=right, fill=white, inner sep=1pt] at (15.8,3.2) {eigener Anschluss,\\nur Realbetrieb};
  % Legende
  \draw[ziel] (1.3,-1.75) -- (2.2,-1.75);
  \node[klein, anchor=west] at (2.3,-1.75) {Ziele und Befehle};
  \draw[zust] (6.0,-1.75) -- (6.9,-1.75);
  \node[klein, anchor=west] at (7.0,-1.75) {Zustände und Meldungen};
\end{tikzpicture}
 in einer figure-Umgebung. Nicht skalieren.
\begin{tikzpicture}[ba]
  \tikzset{
    bw/.style={ba block, text width=31mm, minimum height=10mm},
    fk/.style={ba block hervor, minimum height=17mm, minimum width=0pt},
    sub/.style={ba block, font=\footnotesize, minimum height=10mm, minimum width=0pt, inner sep=2pt},
    ziel/.style={ba pfeil},
    zust/.style={ba pfeil, dash pattern=on 3pt off 2pt},
    klein/.style={font=\footnotesize, text=baGrau},
  }
  % ---------------- Ebene 1: Bedienwege ----------------
  \node[klein, anchor=south west] at (1.3,11.45) {Bedienwege (rechnen keine Bewegung, senden Ziele und lesen Zustände)};
  \node[bw] (B1) at (3.0,10.65) {Bedien-\\oberfläche};
  \node[bw] (B2) at (6.6,10.65) {Gamepad};
  \node[bw] (B3) at (10.2,10.65) {Smartphone};
  \node[bw] (B4) at (13.8,10.65) {Externe\\Programme};
  % Sammelleitung ROS 2
  \draw[ba normal] (3.0,9.95) -- (13.8,9.95);
  \foreach \x in {3.0,6.6,10.2,13.8} {\draw[ba normal] (\x,10.15) -- (\x,9.95);}
  \node[klein, fill=white, inner sep=1.5pt] at (4.3,9.95) {ROS 2};
  % ---------------- Container Umgebung ----------------
  \draw[ba hilfslinie] (1.2,5.25) rectangle (15.6,9.35);
  \node[klein, anchor=south west] at (1.3,5.3) {Container Umgebung};
  % Ebene 2: Funktionsknoten
  \node[fk, text width=22mm] (RZ) at (2.5,7.4) {Roboter-\\zustand};
  \node[fk, text width=28mm] (AR) at (5.9,7.4) {Armregelung\\mit Planung};
  \node[fk, minimum width=51mm, minimum height=24mm, anchor=center] (FB) at (10.55,7.4) {};
  \node[font=\small, anchor=north] at (10.55,8.55) {Fortbewegung};
  \node[sub, text width=21mm] (ER) at (9.2,7.15) {Ersatzregler\\(Simulation)};
  \node[sub, text width=21mm] (OA) at (11.9,7.15) {Anbindung Onboard-Regler (Realbetrieb)};
  \node[fk, text width=19mm] (HA) at (14.35,7.4) {Hände};
  % Pfeile Bedienwege -> Funktionsknoten (Ziele nach unten, Zustände nach oben)
  \draw[ziel] (5.6,9.95) -- (5.6,8.6);
  \draw[zust] (6.2,8.6) -- (6.2,9.95);
  \draw[ziel] (8.6,9.95) -- (8.6,8.6);
  \draw[zust] (12.5,8.6) -- (12.5,9.95);
  \draw[ziel] (14.05,9.95) -- (14.05,8.6);
  \draw[zust] (14.65,8.6) -- (14.65,9.95);
  \draw[zust] (2.5,8.25) -- (2.5,9.95);
  % Schutzmechanismen
  \node[ba block hervor, minimum height=6mm, minimum width=108mm, font=\footnotesize] (SM) at (10.1,5.95)
    {Schutzmechanismen: begrenzen \textperiodcentered{} vermeiden \textperiodcentered{} überwachen \textperiodcentered{} anhalten};
  \draw[ba duenn, color=baGruen] (5.9,6.55) -- (5.9,6.25);
  \draw[ba duenn, color=baGruen] (10.55,6.55) -- (10.55,6.25);
  \draw[ba duenn, color=baGruen] (14.35,6.55) -- (14.35,6.25);
  % ---------------- Ebene 3: Herstellerschnittstelle ----------------
  \node[ba block, fill=baHellgrau!40, minimum height=8mm, minimum width=143mm] (DDS) at (8.4,4.2)
    {Herstellerschnittstelle (Unitree-DDS)};
  % Befehle nach unten, Zustände nach oben zwischen Funktionsknoten und Schnittstelle
  \draw[zust] (2.5,4.6) -- (2.5,6.55);
  \draw[ziel] (5.6,5.65) -- (5.6,4.6);
  \draw[zust] (6.2,4.6) -- (6.2,5.65);
  \draw[ziel] (9.5,5.65) -- (9.5,4.6);
  \draw[zust] (10.1,4.6) -- (10.1,5.65);
  \draw[ziel] (11.5,5.65) -- (11.5,4.6);
  \draw[zust] (12.1,4.6) -- (12.1,5.65);
  \draw[ziel] (14.05,5.65) -- (14.05,4.6);
  \draw[zust] (14.65,4.6) -- (14.65,5.65);
  % ---------------- Umschalter ----------------
  \node[draw=baAkzent, ba normal, fill=white, rounded corners=1.5pt, text=black, align=center,
        minimum height=9mm, inner sep=3pt] (SW) at (8.4,2.75) {Umschalter\\Simulation / Realbetrieb};
  \draw[baAkzent, ba normal] (8.4,3.8) -- (SW.north);
  % ---------------- Ebene 4: Simulation oder Roboter ----------------
  \draw[ba gruppe] (1.3,-1.1) rectangle (7.0,1.5);
  \node[klein, anchor=north west] at (1.4,1.5) {Container Simulation};
  \node[ba block, minimum width=0pt, text width=22mm] (SIM) at (2.9,-0.1) {MuJoCo-Simulation};
  \node[ba block, minimum width=0pt, text width=22mm] (SH) at (5.5,-0.1) {Simulierte Hände};
  \draw[ba gruppe] (9.8,-1.1) rectangle (15.5,1.5);
  \node[klein, anchor=north west] at (9.9,1.5) {Realbetrieb};
  \node[ba block, minimum width=0pt, text width=22mm] (RB) at (11.4,-0.1) {Unitree G1 mit Onboard-Regler};
  \node[ba block, minimum width=0pt, text width=22mm] (RH) at (14.0,-0.1) {Inspire-Hände};
  \draw[baAkzent, ba normal, <->] (SW.west) -| (4.15,1.5);
  \draw[baAkzent, ba normal, <->] (SW.east) -| (12.65,1.5);
  \node[font=\footnotesize, anchor=east, align=right] at (4.05,2.1) {Domain 1,\\Schnittstelle Loopback};
  \node[font=\footnotesize, anchor=west, align=left] at (12.75,2.1) {Domain 0,\\Netzwerkkarte};
  % Hände real: eigener Anschluss neben der Herstellerschnittstelle
  \draw[ba normal, dash pattern=on 1pt off 2pt, <->] (HA.east) -- (15.85,7.4) -- (15.85,-0.1) -- (RH.east);
  \node[klein, anchor=east, align=right, fill=white, inner sep=1pt] at (15.8,3.2) {eigener Anschluss,\\nur Realbetrieb};
  % Legende
  \draw[ziel] (1.3,-1.75) -- (2.2,-1.75);
  \node[klein, anchor=west] at (2.3,-1.75) {Ziele und Befehle};
  \draw[zust] (6.0,-1.75) -- (6.9,-1.75);
  \node[klein, anchor=west] at (7.0,-1.75) {Zustände und Meldungen};
\end{tikzpicture}
 in einer figure-Umgebung. Nicht skalieren.
\begin{tikzpicture}[ba]
  \tikzset{
    bw/.style={ba block, text width=31mm, minimum height=10mm},
    fk/.style={ba block hervor, minimum height=17mm, minimum width=0pt},
    sub/.style={ba block, font=\footnotesize, minimum height=10mm, minimum width=0pt, inner sep=2pt},
    ziel/.style={ba pfeil},
    zust/.style={ba pfeil, dash pattern=on 3pt off 2pt},
    klein/.style={font=\footnotesize, text=baGrau},
  }
  % ---------------- Ebene 1: Bedienwege ----------------
  \node[klein, anchor=south west] at (1.3,11.45) {Bedienwege (rechnen keine Bewegung, senden Ziele und lesen Zustände)};
  \node[bw] (B1) at (3.0,10.65) {Bedien-\\oberfläche};
  \node[bw] (B2) at (6.6,10.65) {Gamepad};
  \node[bw] (B3) at (10.2,10.65) {Smartphone};
  \node[bw] (B4) at (13.8,10.65) {Externe\\Programme};
  % Sammelleitung ROS 2
  \draw[ba normal] (3.0,9.95) -- (13.8,9.95);
  \foreach \x in {3.0,6.6,10.2,13.8} {\draw[ba normal] (\x,10.15) -- (\x,9.95);}
  \node[klein, fill=white, inner sep=1.5pt] at (4.3,9.95) {ROS 2};
  % ---------------- Container Umgebung ----------------
  \draw[ba hilfslinie] (1.2,5.25) rectangle (15.6,9.35);
  \node[klein, anchor=south west] at (1.3,5.3) {Container Umgebung};
  % Ebene 2: Funktionsknoten
  \node[fk, text width=22mm] (RZ) at (2.5,7.4) {Roboter-\\zustand};
  \node[fk, text width=28mm] (AR) at (5.9,7.4) {Armregelung\\mit Planung};
  \node[fk, minimum width=51mm, minimum height=24mm, anchor=center] (FB) at (10.55,7.4) {};
  \node[font=\small, anchor=north] at (10.55,8.55) {Fortbewegung};
  \node[sub, text width=21mm] (ER) at (9.2,7.15) {Ersatzregler\\(Simulation)};
  \node[sub, text width=21mm] (OA) at (11.9,7.15) {Anbindung Onboard-Regler (Realbetrieb)};
  \node[fk, text width=19mm] (HA) at (14.35,7.4) {Hände};
  % Pfeile Bedienwege -> Funktionsknoten (Ziele nach unten, Zustände nach oben)
  \draw[ziel] (5.6,9.95) -- (5.6,8.6);
  \draw[zust] (6.2,8.6) -- (6.2,9.95);
  \draw[ziel] (8.6,9.95) -- (8.6,8.6);
  \draw[zust] (12.5,8.6) -- (12.5,9.95);
  \draw[ziel] (14.05,9.95) -- (14.05,8.6);
  \draw[zust] (14.65,8.6) -- (14.65,9.95);
  \draw[zust] (2.5,8.25) -- (2.5,9.95);
  % Schutzmechanismen
  \node[ba block hervor, minimum height=6mm, minimum width=108mm, font=\footnotesize] (SM) at (10.1,5.95)
    {Schutzmechanismen: begrenzen \textperiodcentered{} vermeiden \textperiodcentered{} überwachen \textperiodcentered{} anhalten};
  \draw[ba duenn, color=baGruen] (5.9,6.55) -- (5.9,6.25);
  \draw[ba duenn, color=baGruen] (10.55,6.55) -- (10.55,6.25);
  \draw[ba duenn, color=baGruen] (14.35,6.55) -- (14.35,6.25);
  % ---------------- Ebene 3: Herstellerschnittstelle ----------------
  \node[ba block, fill=baHellgrau!40, minimum height=8mm, minimum width=143mm] (DDS) at (8.4,4.2)
    {Herstellerschnittstelle (Unitree-DDS)};
  % Befehle nach unten, Zustände nach oben zwischen Funktionsknoten und Schnittstelle
  \draw[zust] (2.5,4.6) -- (2.5,6.55);
  \draw[ziel] (5.6,5.65) -- (5.6,4.6);
  \draw[zust] (6.2,4.6) -- (6.2,5.65);
  \draw[ziel] (9.5,5.65) -- (9.5,4.6);
  \draw[zust] (10.1,4.6) -- (10.1,5.65);
  \draw[ziel] (11.5,5.65) -- (11.5,4.6);
  \draw[zust] (12.1,4.6) -- (12.1,5.65);
  \draw[ziel] (14.05,5.65) -- (14.05,4.6);
  \draw[zust] (14.65,4.6) -- (14.65,5.65);
  % ---------------- Umschalter ----------------
  \node[draw=baAkzent, ba normal, fill=white, rounded corners=1.5pt, text=black, align=center,
        minimum height=9mm, inner sep=3pt] (SW) at (8.4,2.75) {Umschalter\\Simulation / Realbetrieb};
  \draw[baAkzent, ba normal] (8.4,3.8) -- (SW.north);
  % ---------------- Ebene 4: Simulation oder Roboter ----------------
  \draw[ba gruppe] (1.3,-1.1) rectangle (7.0,1.5);
  \node[klein, anchor=north west] at (1.4,1.5) {Container Simulation};
  \node[ba block, minimum width=0pt, text width=22mm] (SIM) at (2.9,-0.1) {MuJoCo-Simulation};
  \node[ba block, minimum width=0pt, text width=22mm] (SH) at (5.5,-0.1) {Simulierte Hände};
  \draw[ba gruppe] (9.8,-1.1) rectangle (15.5,1.5);
  \node[klein, anchor=north west] at (9.9,1.5) {Realbetrieb};
  \node[ba block, minimum width=0pt, text width=22mm] (RB) at (11.4,-0.1) {Unitree G1 mit Onboard-Regler};
  \node[ba block, minimum width=0pt, text width=22mm] (RH) at (14.0,-0.1) {Inspire-Hände};
  \draw[baAkzent, ba normal, <->] (SW.west) -| (4.15,1.5);
  \draw[baAkzent, ba normal, <->] (SW.east) -| (12.65,1.5);
  \node[font=\footnotesize, anchor=east, align=right] at (4.05,2.1) {Domain 1,\\Schnittstelle Loopback};
  \node[font=\footnotesize, anchor=west, align=left] at (12.75,2.1) {Domain 0,\\Netzwerkkarte};
  % Hände real: eigener Anschluss neben der Herstellerschnittstelle
  \draw[ba normal, dash pattern=on 1pt off 2pt, <->] (HA.east) -- (15.85,7.4) -- (15.85,-0.1) -- (RH.east);
  \node[klein, anchor=east, align=right, fill=white, inner sep=1pt] at (15.8,3.2) {eigener Anschluss,\\nur Realbetrieb};
  % Legende
  \draw[ziel] (1.3,-1.75) -- (2.2,-1.75);
  \node[klein, anchor=west] at (2.3,-1.75) {Ziele und Befehle};
  \draw[zust] (6.0,-1.75) -- (6.9,-1.75);
  \node[klein, anchor=west] at (7.0,-1.75) {Zustände und Meldungen};
\end{tikzpicture}
